\expandafter\def\csname HistoricalNoteText559\endcsname{In Beziehung auf diese Einrichtung mag hier noch einiges bemerkt werden. Das Vorhandensein einer Marke zu der erwähnten Prüfung muss als ein unerlässliches Erforderniss für die Zuverlässigkeit der Beobachtungen betrachtet, und also bei der Errichtung eines neuen Gebäudes nothwendig gehörig Bedacht darauf genommen werden. Vor Erbauung des hiesigen magnetischen Observatoriums war auch in Erwägung gekommen, ob es nicht besser sei, diese Marke auf einem eignen besonders fundirten Postament im Innern des Saals anzubringen, als an der von aussen der Witterung ausgesetzten Wand. Man entschied sich für das letztere, da man sonst entweder die Entfernung des Beobachters von der Nadel hätte verringern, oder den Vortheil, Marke und Scale bei einerlei Ocularstellung deutlich zu sehen, aufgeben, oder dem Saale eine noch grössere Länge geben müssen. was auf dem bestimmten Platze nicht einmal thunlich gewesen wäre. Ein künstliches Surrogat anstatt einer Marke anzubringen wurde aus mehrern Gründen für verwerflich gehalten. Auch hielt man die Besorgniss, dass der absolute Ort der Marke, wegen Einflusses der Witterung auf die Wand, einer merklichen Änderung unterworfen sein könne, bei einer soliden Ausführung des Baues und bei der sehr geringen Höhe der Marke über der Grundmauer für wenig erheblich, zumal da man in seiner Gewalt hat, so oft man will, die Winkelmessung zwischen der Marke und dem durch das nördliche Fenster sichtbaren Kirchthurme zu wiederholen. Für die Richtigkeit dieser Ansicht spricht jetzt eine dreijährige Erfahrung.}
\SourceLeaf{552}Magnetisirung der Nadel und von der Intensität des horizontalen Theils der erdmagnetischen Kraft abhängig, mithin streng genommen zu verschiedenen Zeiten nicht ganz dieselbe ist. Eine Anleitung zur scharfen Bestimmung der Schwingungsdauer wird in der Folge gegeben werden [S.\ 374 d.\ B.], für den gegenwärtigen Zweck ist aber eine sehr genaue Kenntniss nicht nöthig, und man kann daher nicht allein die kleinen Veränderungen, denen sie unterworfen ist, ignoriren, sondern man darf sich sogar verstatten, anstatt des genauen Werths die nächste volle Secunde zu substituiren, um dadurch zu bewirken, dass die Augenblicke, wo der Beobachter die unter dem Verticalfaden des Fernrohrs erscheinende Stelle des Scalenbildes scharf zu fixiren hat, immer auf volle Secunden fallen. Dies geschieht von selbst, wenn die dem wahren Werth der Schwingungsdauer am nächsten kommende ganze Zahl eine gerade ist. Ist sie aber ungerade, so hat man, um diese Bequemlichkeit nicht zu verlieren. die Wahl unter folgenden drei Mitteln.

I. Man hält sich dennoch an die nächste gerade Zahl, und darf dies um so mehr, je weniger ihr Unterschied von dem wahren Werth eine halbe Einheit übersteigt, je grösser überhaupt die Schwingungsdauer ist, und je vollkommner man immer die Nadel im beinahe beruhigten Zustande zu erhalten vermag. Die Nadel im magnetischen Observatorium zu Göttingen z.\ B.\ hat gegenwärtig eine Schwingungsdauer von $20''64$; allein obgleich die Zahl 21 hier die nächste ist, so kann man sich doch bei den hier obwaltenden Umständen, wo der Schwingungsbogen selten ein Paar Scalentheile übersteigt, meistens unbedenklich an die bequemere Zahl 20 halten, da sich leicht darthun lässt, dass der \textit{daraus} entspringende Fehler in einem partiellen Resultat nicht den zwanzigsten Theil des Schwingungsbogens, und der Fehler des Endresultats nicht den hundertsten Theil übersteigen kann. Dagegen würde einem Beobachter, dessen Nadel die Schwingungsdauer $10''64$ hätte, zumal wenn er eine gleich vollkommene Beruhigung nicht in seiner Gewalt hätte, zu empfehlen sein, die Zahl 11 und eine der folgenden Abänderungen zu wählen.

II. Man wählt zwar die ungerade Zahl, nimmt aber die Beobachtungsaugenblicke, die nach obiger Formel auf halbe Secunden fallen würden, entweder alle eine halbe Secunde später, oder alle eine halbe früher, was offenbar weiter keinen Unterschied macht, als dass nun auch die sämmtlichen Endresultate nicht für die volle Minute der Uhrzeit, sondern für eine halbe Secunde mehr oder weniger gelten.

\SourceLeaf{553}III. Wenn man das Endresultat nicht, wie in dem oben entwickelten Verfahren, auf eine ungerade, sondern auf eine gerade Anzahl partieller Resultate gründet, so fallen die Beobachtungszeiten von selbst auf volle Secunden, die anstatt der wahren Schwingungsdauer angenommene nächste ganze Zahl möge gerade oder ungerade sein. Soll z.\ B.\ das Endresultat von sechs partiellen abhängen, so sind die Beobachtungszeiten
\[
T-3t,\quad T-2t,\quad T-t,\quad T,\quad T+t,\quad T+2t,\quad T+3t
\]

Dies Verfahren, wobei der Einfluss des von der Schwingungsdauer weggelassenen Bruchs im Endresultat noch vollkommner eliminirt wird, als in dem vorhin beschriebenen, ist vorzüglich solchen Beobachtern zu empfehlen, die kleinere Apparate oder Nadeln von vergleichungsweise kurzer Schwingungsdauer gebrauchen.

Es mag noch bemerkt werden, dass, da durch Auflegung eines kleinen Gewichts die Schwingungsdauer der Nadel vergrössert wird, man durch eine schickliche Wahl des Gewichts und der Auflegungsstelle im Stande ist, die Schwingungsdauer äusserst nahe auf eine ganze Zahl von Secunden zu bringen. Dieser Ausweg ist wohl von einigen Beobachtern gewählt, die nicht genug in ihrer Gewalt hatten. etwas grössere Schwingungsbewegungen von ihrer Nadel abzuwehren. Immer aber bleibt dies ein sehr ungenügender Nothbehelf; denn wenn auch unter solchen Umständen das obige Theorem als ganz scharf gilt, so werden doch die Resultate immer einen viel geringern Grad von Genauigkeit haben, weil es unmöglich ist, wenn die Nadel in einer stark augenfälligen Bewegung begriffen ist, den einem bestimmten Secundenschlage entsprechenden Scalentheil, und dessen Bruchtheil, mit derselben Schärfe zu fixiren, als wenn die Langsamkeit der Bewegung eine Veränderung in einer Secunde kaum bemerken lässt. Die Nadel immer gehörig beruhigt zu halten ist daher eine Vorschrift, deren Wichtigkeit nicht genug eingeschärft werden kann.

Gerade dieser Ursache wegen ist es wichtig, dass immer zwischen zwei auf einander folgenden Beobachtungssätzen, nöthigenfalls gehörige Zeit zu einer Beruhigung bleibe. Bei der Nadel des Göttinger magnetischen Observatoriums ist diese Zwischenzeit, unter Anwendung der ersten Methode $3'\,20''$, bei Anwendung der zweiten würde sie $2'\,54''$ sein: in beiden Fällen für Geübte zu obigem Zweck hinreichend. Gewöhnlich benutzen die Beobachter die Zwischenzeit (da
\SourceLeaf{554}das Bedürfniss, zu beruhigen, äusserst selten eintritt), dazu, eine Reinschrift der Beobachtung zu machen, und das Endresultat zu berechnen. Wo hingegen die Nadel eine viel längere Schwingungsdauer hat, und mithin jene Zwischenzeit zwischen zwei Beobachtungssätzen viel kürzer ausfällt, wird eine diesen Übelstand beseitigende Abänderung der obigen Methoden vorzuziehen sein.

Die Abänderung besteht darin, dass man die einzelnen Beobachtungszeiten nicht um eine Schwingungsdauer, sondern um einen aliquoten Theil derselben (die Hälfte, oder den dritten Theil) von einander abstehen lässt. Ausser dem Vortheile, die Aufzeichnungen zu jedem Satz in kürzerer Zeit abzuthun, und also grössere Zwischenzeit zwischen zwei Sätzen zu gewinnen, entgeht man dabei auch der Unannehmlichkeit des erstern Verfahrens, den grössern Theil der Zwischenzeit zwischen zwei Aufzeichnungen unbeschäftigt zu sein. Geübtere Beobachter wenden daher gern das abgeänderte Verfahren selbst da an, wo die Schwingungsdauer nicht eben sehr lang ist. Bei der hiesigen Anstalt machen mehrere Beobachter ihre Aufzeichnungen in Zwischenzeiten von $10''$ (als Hälfte von $20''$), ja selbst von $7''$ (als drittem Theil von $21''$). Einige Beispiele werden das weiter dabei zu Bemerkende am besten erläutern.

\Needspace{15\baselineskip}
\begin{center}\textit{Beobachtung am} 17. \textit{August} 1836, \textit{für} $10^{\mathrm h}\,20'$ \textit{durch Herrn Prof. Ulrich.}\end{center}
\[
\begin{array}{rrr|r|}
10^{\mathrm h}&19'&30''&869.9\\
&&40&871.3\\
&&50&871.7\\
&20&0&870.8\\
&&10&871.0\\
&&20&872.4\\
&&30&872.9
\end{array}\quad
\left.\begin{array}{r}870.80\\871.05\\871.35\\871.60\\871.95\end{array}\right\}871.35
\]
Die zweite Columne enthält die einzelnen Aufzeichnungen, die dritte die partiellen Resultate, und zwar ist $870.80$ das Mittel der ersten und dritten Aufzeichnung und gilt also für $10^{\mathrm h}\,19'\,40''$ u.\ s.\ f.\ Man sieht in diesem aus einer Zeit schneller Veränderung der Declination gewählten Beispiele mit Vergnügen, wie ein geübter Beobachter die Veränderungen in 10 Secunden mit Sicherheit erkennen kann.

\SourceLeaf{555}\Needspace{16\baselineskip}
\begin{center}\textit{Beobachtung am} 25. \textit{März} 1837, \textit{für} $0^{\mathrm h}\,5'$ \textit{durch Herrn Dr. Goldschmidt.}\end{center}
\[
\begin{array}{rrr|r|}
0^{\mathrm h}&4'&32''&847.3\\
&&39&847.2\\
&&46&847.8\\
&&53&848.7\\
&5&0&848.9\\
&&7&848.1\\
&&14&847.0\\
&&21&846.9\\
&&28&847.3
\end{array}\quad
\left.\begin{array}{r}848.00\\848.05\\847.95\\847.85\\847.90\\847.70\end{array}\right\}847.91
\]
Das erste partielle Resultat entspringt hier aus der Combination der ersten und vierten Aufzeichnung, das zweite aus der zweiten und fünften u.\ s.\ w.

In diesen Beispielen war das Submultiplum der zum Grunde gelegten genäherten Schwingungsdauer eine ganze Zahl; wo dies nicht der Fall ist, muss man die Schwingungsdauer in ungleiche Theile zerlegen, was aber keinen Nachtheil hat, wenn man nur die Einrichtung so macht, dass den zu combinirenden Aufzeichnungen immer \textit{derselbe} genäherte Werth der Schwingungsdauer als Zwischenzeit entspreche, und, falls man es der Mühe werth hält, im Protocollauszuge die Zeit, welcher das Endresultat entspricht, mit ihrem Bruchtheile bemerkt. So werden z.\ B.\ die Beobachtungen mit dem 25pfündigen Stabe in der Sternwarte, dessen Schwingungsdauer jetzt $43''14$ ist, wenn man den dafür zu setzenden genäherten Werth $43''$ in vier Theile abtheilen, und das Endresultat auf fünf partielle Resultate gründen will, nach folgendem Schema angestellt:
\[
\begin{array}{rrr|}
0^{\mathrm h}&4'&17''\\
&&28\\
&&39\\
&&49\\
&5&0\\
&&11\\
&&22\\
&&32\\
&&43
\end{array}\quad
\left.\begin{array}{rrr}0^{\mathrm h}&4'&38''5\\&&49.5\\&5&0.5\\&&10.5\\&&21.5\end{array}\right\}0^{\mathrm h}\,5'\,0''1
\]
Hier enthält die erste Columne die Aufzeichnungszeiten, die zweite die Zeiten, für welche die partiellen Resultate eigentlich gelten, und wo natürlich
\SourceLeaf{556}ganz gleichgültig ist, dass das Endresultat genau genommen auf $0^{\mathrm h}\,5'\,0''1$ fällt. Soll das Endresultat auf sechs partielle Resultate gegründet sein, so wird folgendes Schema befolgt:
\[
\begin{array}{rrr|}
0^{\mathrm h}&4'&12'\\
&&22\\
&&33\\
&&44\\
&&55\\
&5&5\\
&&16\\
&&27\\
&&38\\
&&48
\end{array}\quad
\left.\begin{array}{rrr}0^{\mathrm h}&4'&33''5\\&&43.5\\&&54.5\\&5&5.5\\&&16.5\\&&26.5\end{array}\right\}0^{\mathrm h}\,5'\,0''
\]

Am klarsten tritt der Vortheil der abgeänderten Beobachtungsart hervor, wenn der Gang der magnetischen Declination in engern Zwischenzeiten als von fünf zu fünf Minuten verfolgt werden soll. Diese Zwischenzeiten, ausreichend bei dem gewöhnlichen Hergange der Declinationsveränderungen, sind in der That noch zu gross, um den stärkern und schneller wechselnden Änderungen ganz ihr Recht wiederfahren zu lassen, und gerade diese Rücksicht hatte, weil engere Intervalle nicht wohl zur allgemeinen und durchgängigen Regel für die vierundzwanzigstündigen Termine gemacht werden konnten, die Festsetzung der Nebentermine veranlasst, in welchen jedesmal zwei Stunden von drei zu drei Minuten beobachtet werden sollte. Da indessen die Abhaltung dieser Nebentermine an manchen Orten Schwierigkeiten gefunden, und es sich auch so gefügt hat, dass bisher nur in wenigen beträchtliche Bewegungen vorgekommen sind, so ist beschlossen, sie von jetzt an fallen zu lassen, zumal da derselbe wichtige Zweck auch auf andere Art, und selbst noch besser sich wird erreichen lassen. Die Zwischenzeiten von fünf zu fünf Minuten bleiben nach wie vor die Regel; so oft aber das Vorhandensein schneller Declinationsänderung bemerkt wird, werden die Sätze, so lange es als nöthig erscheint, von $2\frac12$ zu $2\frac12$ Minuten ausgeführt. Nach dem, was oben entwickelt ist, wird, anstatt aller weitern Erläuterung, genügen, wenn dem obigen Beispiele vom 17. August [S.\ 548] noch die unmittelbar darauf folgende Beobachtung beigefügt wird:
\SourceLeaf{557}
\[
\begin{array}{rrr|r|}
10^{\mathrm h}&22'&0''&875.0\\
&&10&874.8\\
&&20&876.0\\
&&30&877.1\\
&&40&876.8\\
&&50&876.1\\
&23&0&877.1
\end{array}\quad
\left.\begin{array}{r}875.50\\875.95\\876.40\\876.60\\876.90\end{array}\right\}876.27\quad\text{für }10^{\mathrm h}\,22'\,30''
\]

Die sämmtlichen auswärtigen Theilnehmer werden aufgefordert, es in vorkommenden Fällen auf dieselbe Weise zu halten: es lässt sich nicht zweifeln, dass dann immer für alle grössern Bewegungen eine Menge im engen Detail correspondirender Beobachtungen zusammenkommen, und über die Verhältnisse dieser merkwürdigen Erscheinungen interessante Aufschlüsse geben werden.

Für den Fall, wo man sich beim Beobachten nicht einer Secundenpendeluhr, sondern einer Uhr bedient, die andere Zeittheile schlägt, wird eine besondere Anweisung nicht nöthig sein. Man zählt dann, anstatt der Secunden, die Uhrschläge, und ordnet das Geschäft auf ganz analoge Weise so an, dass alle Beobachtungen auf bestimmte Schläge gemacht werden. Es wird aber eine etwas grössere Aufmerksamkeit erfordert, die Schläge eines Chronometers immer richtig zu zählen, als die Schläge einer Pendeluhr, zumal wenn bei jenem der Zeiger einige Excentricität hat, und deswegen nicht an allen Stellen des Zifferblatts, wo er sollte, genau auf die Secundenstriche springt.

Einige allgemeine Vorsichtsmaassregeln, obwohl zum Theil scheinbare Geringfügigkeiten betreffend, verdienen noch hier erwähnt zu werden, da mancher angehende Beobachter, ohne im voraus aufmerksam darauf gemacht zu sein, sie anfangs leicht übersehen könnte.

Das allererste Erforderniss ist, dass die Nadel völlig frei schwingen könne. Solche Hindernisse der freien Bewegung, die sogleich offenbar ins Auge fallen, wird natürlich jeder Beobachter von selbst wegzuräumen wissen: es gibt aber auch andere, dem Auge sich fast entziehende, die gleichwohl die Beobachtungen ganz verderben können.

In der wärmern Jahrszeit findet sich zuweilen wohl eine Spinne im Kasten ein (am leichtesten, wenn die Seitenöffnung vor dem Spiegel stets offen bleibt), knüpft ein Gewebe oder einen einzelnen Faden zwischen dem Magnetstabe oder dessen Zubehör, und dem Kasten, und hemmt dadurch die freie Bewegung des
\SourceLeaf{558}Magnetstabes. Man thut daher wohl, sich kurz vor jedem Termin erst zu überzeugen, dass der Kasten innen rein ist. Ist der Deckel des Kastens verglaset, so erkennt man die Gegenwart grösserer Insecten oder Gewebe schon von aussen; allein man unterlasse nicht, den Deckel abzuheben und genauer nachzusehen; ja man beruhige sich nicht dabei, wenn man gar keinen Faden sieht, denn in der That reicht, wie öftere Erfahrungen bewiesen haben, auch der allerfeinste dem blossen Auge unsichtbare oder nur unter ganz besonderer Beleuchtung erkennbare Faden schon hin, die freie Bewegung zu hemmen, und die Beobachtungen zu verderben. Um sich gegen solchen, weil unsichtbar gefährlichsten, Feind zu sichern, umfahre man den Magnetstab mit dem Finger, einem Stäbchen oder dergl.\ auf allen Seiten, rechts, links, vorne, hinten, oben und unten, wodurch ein solcher Faden, wenn einer da war, zerrissen wird. Fast eben so sicher erreicht man dieselbe Wirkung dadurch, dass man den Stab in sehr grosse Schwingungen versetzt. Es verdient noch bemerkt zu werden, dass solche und ähnliche Hindernisse der freien Bewegung allemal mit einer Verminderung der Schwingungsdauer der Nadel verbunden sind, und zwar bewirken selbst äusserst zarte Spinnefäden schon eine sehr bedeutende Verminderung der Schwingungsdauer, wovon unten ein merkwürdiges Beispiel vorkommen wird.

Für die nächtlichen Beobachtungen ist es nothwendig, die Scale zu erleuchten, was in Göttingen in den Terminsbeobachtungen durch zwei \textsc{Argand}sche Lampen geschieht. Über der Lichtflamme findet immer ein Aufströmen erwärmter Luft Statt, und wenn dabei eine der Lampen nahe eben unter dem Fernrohr steht, so hat solche Luftströmung vor dem Objectiv auf die Deutlichkeit des Sehens einen nachtheiligen Einfluss; die Theilstriche der Scale erscheinen zitternd oder wallend. Dieser Übelstand trat in Göttingen bei den ersten Beobachtungen öfters ein, hat aber vollkommen aufgehört, seitdem jede Lampe mit einem seitwärts gebogenen Schorstein aus Kupferblech versehen ist.

Da in die Arbeit zu den Terminsbeobachtungen sich immer eine grössere oder kleinere Zahl von Personen theilen muss, so wird gewöhnlich eine beträchtliche Ungleichheit der Sehweite bei denselben Statt finden: das vollkommen deutliche Sehen ist aber ein durchaus wesentliches Erforderniss für gute Beobachtungen. Wird ein Weitsichtiger, für dessen Auge das Fernrohr zum vollkommen deutlichen Sehen gestellt war, von einem Kurzsichtigen abgelöst, so würde dieser ohne eine Veränderung am Fernrohr gar keine brauchbaren Beobachtungen
\SourceLeaf{559}anstellen können. Die Zuziehung eines Hohlglases würde unbequem und auch wegen des bedeutenden Lichtverlustes nicht anzurathen sein. Das blosse Einschieben der Ocularröhre reicht nicht hin, weil, wenn gleich dadurch das Scalenbild zur Deutlichkeit gebracht wird, doch das Fadenkreuz undeutlich bleiben und gegen das Bild des Gegenstandes eine Parallaxe erhalten würde. Es müsste daher (bei der Einrichtung, die die zu solchen Beobachtungen angewandten Fernröhre zu haben pflegen) zugleich die das Fadenkreuz tragende innere Hülse in der Ocularröhre verschoben und dem Ocularglase näher gebracht werden, was aber eine geübte Hand erfordert, Zeitaufwand veranlasst, und auch aus andern Gründen für den vorliegenden Fall nicht zu empfehlen ist. Man kann aber dem Bedürfniss auf eine sehr einfache Art abhelfen, wenn man sich folgendes Verfahren zur Regel macht. Die Ocularröhre im Fernrohr und das Fadenkreuz in derselben ist vor Anfang der Beobachtungen so gestellt, dass der Kurzsichtigste unter den Beobachtern Fadenkreuz und Scalenbild zugleich vollkommen deutlich sieht: so oft ein weitsichtiger Beobachter an die Reihe kommt, hat derselbe, ohne die Ocularröhre oder das Fadenkreuz in derselben zu verrücken, nur das dem Auge nächste Glas so weit zurückzuschrauben, dass er das Fadenkreuz vollkommen scharf sieht, womit denn ein völlig deutliches Sehen des Scalenbildes schon von selbst verbunden ist. Ein später eintretender Kurzsichtiger hat dann nur dieses Glas so viel sein Auge erfordert wieder hineinzuschrauben.

Zur Prüfung des unverrückten Standes des Fernrohrs dient eine Marke, die in solcher Entfernung angebracht ist, dass sie bei der zum deutlichen Sehen des Scalenbildes erforderlichen Ocularstellung gleichfalls deutlich erscheint, und im Göttinger magnetischen Observatorium blos in einem feinen verticalen Strich an der nördlichen Wand besteht\HistoricalNote{559}.

\SourceLeaf{560}Vor Anfang der Beobachtungen hat man das Fernrohr nach der Marke zu richten, nachher von Zeit zu Zeit die Prüfung zu wiederholen, und sobald sich eine Abweichung zeigt, die optische Axe des Fernrohrs wieder in die vorige Verticalebene zurückzubringen. Hat man neben der Marke auf beiden Seiten noch eine Eintheilung angebracht, so erkennt man dadurch zugleich die Grösse der nöthig gewordenen Correction, wobei jedoch erinnert werden mag, dass jene Theile, wenn sie auch wie die Scalentheile Millimeter sind, genau genommen nicht ganz denselben Werth in Secunden haben werden, wie die letztern. In Ermangelung jener Hülfseintheilung kann man sich jedoch schon begnügen, die Grösse der gefundenen Abweichung in Scalentheilen, blos wie diese erscheinen, nach dem Augenmaass zu schätzen.

Die Beobachtungen werden am verticalen Faden des Fadenkreuzes gemacht, während der horizontale blos dient, ungefähr die Mitte des erstern zu bezeichnen. Damit es keinen Unterschied mache, ob man die Theile der Scale etwas höher oder tiefer im Gesichtsfelde erscheinen lasse, muss das Fadenkreuz eine solche Stellung haben, dass ein festes auf der Kreuzung der Fäden sich abbildendes Object genau auf dem Verticalfaden bleibe, wenn man das Fernrohr etwas auf und nieder bewegt. Auch zu dieser Berichtigung, die übrigens selten wiederholt zu werden braucht, wenn man die Stellung der Ocularröhre unverändert lässt, dient die Marke.

Der von der Mitte des Objectivs herabhängende Lothfaden ist der Scale so nahe, dass das Bild von beiden im Fernrohr mit gleicher Deutlichkeit erscheint, und man also den Theilstrich, welchen jener Faden deckt, sehr scharf beobachten kann. Man bringt die Scale so an, dass jener Punkt der Scale ihre Mitte, oder ein willkürlich dafür angenommener Theilstrich ist. Die Prüfung des unverrückten Standes der Scale ist im Laufe der Beobachtungen von Zeit zu Zeit zu wiederholen; es ist jedoch nicht nöthig, wenn man eine kleine Änderung findet, die Scale wieder in die vorige Stellung zu bringen, sondern es reicht hin, den dem Lothfaden entsprechenden Theilungspunkt im Protocoll zu bemerken.

\SourceLeaf{561}Hiebei ist es jedoch vielleicht nicht überflüssig, auf ein Paar Kleinigkeiten besonders aufmerksam zu machen.

Es wird zwar vorausgesetzt, dass Magnetometer und Fernrohr so aufgestellt sind, dass der mittlere Stand der magnetischen Declination ungefähr der Mitte der Scale entspricht. Allein zur Zeit beträchtlicher Variationen kommt nicht selten diese Mitte ganz aus dem Gesichtsfelde, und man kann so obige Prüfung nicht vornehmen. Hat man zu solcher Zeit Veranlassung zu jener Prüfung, so muss man den Beruhigungsmagnet einen seinem gewöhnlichen Gebrauch gerade entgegengesetzten Dienst leisten lassen, nemlich die Nadel des Magnetometers in solche Schwingungen versetzen, die bis zu der gesuchten Stelle oder ein wenig darüber hinausgehen, wodurch man also Gelegenheit erhält, den Lothfaden in der Mitte des Gesichtsfeldes zu sehen, und zwar in einer solchen Zeit einer Schwingungsperiode, wo die Geschwindigkeit der Bewegung gering, also das scharfe Auffassen des entsprechenden Theilungspunkts nicht gehindert ist. Da man, wenn dergleichen im Laufe der Beobachtungen vorfällt, sogleich wieder zur Beruhigung schreiten muss, um wo möglich den folgenden Beobachtungssatz nicht zu verlieren, so erhellt, wie nützlich es ist, mit dem Gebrauch des Beruhigungsmagnets recht vertraut zu sein.

Im umgekehrten Fall, nemlich so oft die Declination in die Nähe der Mitte der Scale trifft, ist für Ungeübte eine andere Warnung nöthig, nemlich den Lothfaden nicht mit dem Verticalfaden des Fernrohrs zu verwechseln. Am hiesigen Apparat erscheinen in der That beide einander so sehr gleich, dass bei sehr ruhigem Stande der Nadel ohne ein Paar an letzterem Faden haftende Stäbchen eine Verwechslung wohl möglich wäre, und an einem andern Orte ist wirklich früher einmal der Fall vorgekommen, dass ein Beobachter eine halbe Stunde hindurch die Nadel völlig stationär fand, während er immer den unrechten Faden beobachtet hatte. Da bei einer sehr grossen Annäherung beider Fäden das Beobachten immer ein wenig erschwert wird, so thut man wohl, in einem solchen Falle den Lothfaden eine Zeitlang zu beseitigen.

Was die Form der Mittheilung betrifft, so pflegen einige die Beobachtungen ganz \textit{in extenso}, andere die partiellen und die Endresultate, und mehrere blos die letztern einzusenden. In der Voraussetzung, dass vorher die Rechnungen durchgesehen und die mitgetheilten Zahlen collationirt sind, kann dieser Auszug auch genügen: indessen werden die Beobachtungen selbst, um erforderlichen
